\documentclass[11pt,a4paper]{article}

\usepackage[utf8]{inputenc}
\usepackage[T1]{fontenc}
\usepackage[french]{babel}
\usepackage{amsmath,amssymb}
\usepackage{geometry}
\geometry{margin=2.5cm}
\usepackage{tikz}
\usetikzlibrary{arrows.meta,decorations.markings}
\usepackage{array}
\usepackage{enumitem}

\setlength{\parindent}{0pt}

% --- Réglage global de la hauteur des lignes de tableaux ---
\renewcommand{\arraystretch}{2.2}
\setlength{\extrarowheight}{2pt}

\title{Trigonométrie}
\author{}
\date{\today}

\begin{document}

\maketitle

% =====================================================================
\clearpage
\section*{Partie 1 — Conversion degré vers radian et vice-versa}
% =====================================================================

\[
180^\circ=\pi\ \text{radians}.
\]
Donc :
\[
\text{radians}=\text{degrés}\times \dfrac{\pi}{180},
\qquad
\text{degrés}=\text{radians}\times \dfrac{180}{\pi}.
\]

Exemples :
\[
30^\circ=\dfrac{\pi}{6},\qquad
45^\circ=\dfrac{\pi}{4},\qquad
60^\circ=\dfrac{\pi}{3},\qquad
90^\circ=\dfrac{\pi}{2},\qquad
180^\circ=\pi.
\]
\[
\dfrac{2\pi}{3}=120^\circ,\qquad
\dfrac{3\pi}{4}=135^\circ,\qquad
\dfrac{5\pi}{6}=150^\circ.
\]

% =====================================================================
\clearpage
\section*{Partie 2 — Avant-propos : que signifie « modulo $2\pi$ » ?}
% =====================================================================

\subsection*{2.1. Le problème du « tour de trop »}

Sur le cercle trigonométrique, un angle est repéré par la position du point $M$. Or plusieurs nombres peuvent repérer \emph{le même} point $M$ : il suffit d'ajouter ou de retirer un ou plusieurs tours complets.

\begin{center}
\begin{tikzpicture}[scale=2]
  \draw[->] (-1.4,0) -- (1.4,0) node[right] {$\vec{i}$};
  \draw[->] (0,-1.4) -- (0,1.4) node[above] {$\vec{j}$};
  \draw (0,0) circle (1);

  \fill[red] (30:1) circle (2pt);
  \node[red,above right] at (30:1) {$M$};

  \draw[->,red,thick] (0.4,0) arc (0:30:0.4);
  \node[red] at (15:0.55) {\small $\dfrac{\pi}{6}$};

  \draw[->,blue,thick] (0.65,0) arc (0:390:0.65);
  \node[blue] at (-0.85,-0.65) {\small $\dfrac{\pi}{6}+2\pi$};

  \draw[->,green!60!black,thick] (0.9,0) arc (0:-330:0.9);
  \node[green!60!black] at (0.8,0.85) {\small $\dfrac{\pi}{6}-2\pi$};
\end{tikzpicture}
\end{center}

\textbf{Constations :}
\begin{itemize}
  \item $\dfrac{\pi}{6}$, \quad $\dfrac{\pi}{6}+2\pi$, \quad $\dfrac{\pi}{6}-2\pi$ \quad désignent \textbf{le même point $M$} sur le cercle.
  \item En radians, un tour complet vaut $2\pi$. En degrés, un tour complet vaut $360^\circ$.
  \item Ajouter $2\pi$ (ou $360^\circ$) revient donc à « faire un tour de plus » et à revenir au même endroit.
\end{itemize}

\subsection*{2.2. Un même point, une infinité de nombres}

On peut ajouter ou retirer autant de tours que l'on veut. Pour un même point $M$, il existe une \textbf{infinité} de nombres qui le repèrent :
\[
\theta,\quad \theta+2\pi,\quad \theta-2\pi,\quad \theta+4\pi,\quad \theta-6\pi,\quad\ldots
\]
Plus généralement :
\[
\theta+2k\pi \quad \text{avec } k\in\mathbb{Z}.
\]

\begin{center}
\begin{tikzpicture}[scale=1.8]
  \draw[->] (-1.4,0) -- (1.4,0) node[right] {$\vec{i}$};
  \draw[->] (0,-1.4) -- (0,1.4) node[above] {$\vec{j}$};
  \draw (0,0) circle (1);

  \fill[red] (60:1) circle (2pt);
  \node[red,above right] at (60:1) {$M$};

  \draw[->,red,thick] (0.35,0) arc (0:60:0.35);
  \draw[->,blue,thick] (0.55,0) arc (0:420:0.55);
  \draw[->,green!60!black,thick] (0.75,0) arc (0:-300:0.75);
  \draw[->,orange,thick] (0.95,0) arc (0:780:0.95);
\end{tikzpicture}
\end{center}

Sur cette figure, quatre arcs différents partent tous de l'axe des abscisses et aboutissent au même point $M$ (d'angle $\dfrac{\pi}{3}$) :
\[
\dfrac{\pi}{3},\qquad
\dfrac{\pi}{3}+2\pi,\qquad
\dfrac{\pi}{3}-2\pi,\qquad
\dfrac{\pi}{3}+4\pi.
\]

\subsection*{2.3. La notation « mod $2\pi$ »}

Pour dire que deux nombres repèrent le même point, on écrit :
\[
\theta \equiv \theta' \pmod{2\pi}.
\]
Cela se lit : « $\theta$ est congru à $\theta'$ modulo $2\pi$ ». Cela signifie :
\[
\text{il existe } k\in\mathbb{Z} \text{ tel que } \theta=\theta'+2k\pi.
\]

\paragraph{Exemples.}
\[
\dfrac{\pi}{6}\equiv \dfrac{13\pi}{6}\pmod{2\pi}
\quad\text{car}\quad
\dfrac{13\pi}{6}=\dfrac{\pi}{6}+2\pi.
\]
\[
\dfrac{\pi}{3}\equiv -\dfrac{5\pi}{3}\pmod{2\pi}
\quad\text{car}\quad
-\dfrac{5\pi}{3}=\dfrac{\pi}{3}-2\pi.
\]
\[
\pi\equiv 3\pi\equiv -5\pi \pmod{2\pi}.
\]

\subsection*{2.4. Comparaison avec les degrés : « mod $360^\circ$ »}

En degrés, le même phénomène s'écrit avec $360^\circ$ à la place de $2\pi$ :
\[
\theta \equiv \theta' \pmod{360^\circ}
\iff
\text{il existe } k\in\mathbb{Z} \text{ tel que } \theta=\theta'+360k.
\]

\begin{center}
\begin{tabular}{|c|c|c|c|}
\hline
\textbf{En degrés} & \textbf{En radians} & \textbf{Même point ?} & \textbf{Différence}\\
\hline
$30^\circ$ & $\dfrac{\pi}{6}$ & oui & référence\\
\hline
$390^\circ$ & $\dfrac{\pi}{6}+2\pi=\dfrac{13\pi}{6}$ & oui & $+360^\circ$\\
\hline
$-330^\circ$ & $\dfrac{\pi}{6}-2\pi=-\dfrac{11\pi}{6}$ & oui & $-360^\circ$\\
\hline
$750^\circ$ & $\dfrac{\pi}{6}+4\pi=\dfrac{25\pi}{6}$ & oui & $+720^\circ$\\
\hline
$150^\circ$ & $\dfrac{5\pi}{6}$ & non & $+120^\circ$\\
\hline
\end{tabular}
\end{center}

\begin{center}
\begin{tikzpicture}[scale=1.6]
  \draw[->] (-1.3,0) -- (1.3,0) node[right] {$\vec{i}$};
  \draw[->] (0,-1.3) -- (0,1.3) node[above] {$\vec{j}$};
  \draw (0,0) circle (1);

  \fill[red] (30:1) circle (2.5pt);
  \node[red,above right] at (30:1) {\small $30^\circ$};

  \fill[blue] (150:1) circle (2.5pt);
  \node[blue,above left] at (150:1) {\small $150^\circ$};

  \draw[->,red,thick] (0.3,0) arc (0:30:0.3);
  \draw[->,blue,thick] (0.5,0) arc (0:150:0.5);
\end{tikzpicture}
\end{center}

\subsection*{2.5. Pourquoi c'est indispensable en trigonométrie ?}

\begin{itemize}
  \item \textbf{Pour repérer un point} $M$ du cercle, un seul nombre suffit ; on choisit la \emph{mesure principale}, comprise entre $-\pi$ et $\pi$ (soit entre $-180^\circ$ et $180^\circ$).
  \item \textbf{Deux nombres congrus mod $2\pi$ donnent le même point}, donc le même cosinus et le même sinus.
  \item \textbf{Pour comparer deux angles}, on les ramène d'abord dans le même « tour » en utilisant la congruence modulo $2\pi$.
\end{itemize}

\paragraph{Conséquence immédiate.}
Si $\theta\equiv\theta'\pmod{2\pi}$, alors
\[
\cos\theta=\cos\theta'
\qquad\text{et}\qquad
\sin\theta=\sin\theta'.
\]
C'est ce que l'on utilisera dans toute la suite du cours.

% =====================================================================
\clearpage
\section*{Partie 3 — Mesure principale en degré et en radian}
% =====================================================================

La mesure principale d'un angle orienté est l'unique mesure $\alpha$ telle que
\[
-\pi<\alpha\le \pi
\]
en radians, ou de façon équivalente
\[
-180^\circ<\alpha\le 180^\circ
\]
en degrés. Toute autre mesure s'écrit $\alpha+2k\pi$ avec $k\in\mathbb{Z}$.

Autrement dit : parmi tous les nombres qui repèrent un même point $M$ du cercle, la mesure principale est \textbf{le seul} qui appartienne à l'intervalle $]-\pi\,;\,\pi]$ (ou $]-180^\circ\,;\,180^\circ]$).

\paragraph{Exemples.}
\[
\dfrac{7\pi}{3}\equiv \dfrac{7\pi}{3}-2\pi=\dfrac{\pi}{3}\pmod{2\pi}.
\]
\[
-\dfrac{5\pi}{4}\equiv -\dfrac{5\pi}{4}+2\pi=\dfrac{3\pi}{4}\pmod{2\pi}.
\]
\[
750^\circ\equiv 750^\circ-2\times 360^\circ=30^\circ\pmod{360^\circ}.
\]
\[
-200^\circ\equiv -200^\circ+360^\circ=160^\circ\pmod{360^\circ}.
\]

% =====================================================================
\clearpage
\section*{Partie 4 — Symétries et impacts : tableau récapitulatif}
% =====================================================================

Soit $(O,\vec{i},\vec{j})$ un repère orthonormé direct. Pour tout point $M(x,y)$, on note
\[
\theta=\left(\vec{i},\overrightarrow{OM}\right)\pmod{2\pi}.
\]

\begin{center}
\begin{tabular}{|c|c|c|}
\hline
\textbf{Symétrie} & \textbf{Coordonnées de $M'$} & \textbf{Angle image} \\
\hline
Axe $(O,\vec{i})$ & $(x,-y)$ & $-\theta$ \\
\hline
Axe $(O,\vec{j})$ & $(-x,y)$ & $\pi-\theta$ \\
\hline
Centre $O$ & $(-x,-y)$ & $\pi+\theta$ \\
\hline
\end{tabular}
\end{center}

% =====================================================================
\clearpage
\section*{Partie 5 — Rotation d'angle $\pm 90^\circ$}
% =====================================================================

\subsection*{5.1. Impact de ces rotations sur les coordonnées de l'image d'un point $M$}

Soit $M(x,y)$ un point du plan.

\paragraph{Rotation d'angle $+90^\circ$ (sens direct).}

L'image $M'$ de $M$ par la rotation de centre $O$ et d'angle $+90^\circ$ a pour coordonnées
\[
M'(x',y')=(-y\,;\,x).
\]
\textit{Justification géométrique :} la rotation de $+90^\circ$ fait passer l'axe des abscisses sur l'axe des ordonnées. L'abscisse $x$ de $M$ devient donc l'ordonnée de $M'$, et l'ordonnée $y$ de $M$ devient l'opposé de l'abscisse de $M'$.

\paragraph{Rotation d'angle $-90^\circ$ (sens indirect).}

L'image $M'$ de $M$ par la rotation de centre $O$ et d'angle $-90^\circ$ a pour coordonnées
\[
M'(x',y')=(y\,;\,-x).
\]
\textit{Justification géométrique :} c'est la rotation inverse de la précédente.

\subsection*{5.2. Impact de ces rotations sur l'image de l'angle orienté $(\vec{i},\overrightarrow{OM})$}

Soit $\theta=\left(\vec{i},\overrightarrow{OM}\right)$.

\begin{itemize}
  \item Rotation d'angle $+90^\circ=+\dfrac{\pi}{2}$ :
  \[
  \theta'=\theta+\dfrac{\pi}{2}\pmod{2\pi}.
  \]

  \item Rotation d'angle $-90^\circ=-\dfrac{\pi}{2}$ :
  \[
  \theta'=\theta-\dfrac{\pi}{2}\pmod{2\pi}.
  \]
\end{itemize}

% =====================================================================
\clearpage
\section*{Partie 6 — Cercle trigonométrique}
% =====================================================================

Le cercle trigonométrique est le cercle de centre $O$ et de rayon $1$, orienté dans le sens direct. Son équation est
\[
x^2+y^2=1.
\]
À tout réel $\theta$ correspond un unique point $M$ du cercle tel que
\[
\left(\vec{i},\overrightarrow{OM}\right)\equiv \theta\pmod{2\pi}.
\]

\begin{center}
\begin{tikzpicture}[scale=2.2]
  \draw[->] (-1.25,0) -- (1.25,0) node[right] {$\vec{i}$};
  \draw[->] (0,-1.25) -- (0,1.25) node[above] {$\vec{j}$};
  \draw (0,0) circle (1);
  \draw[->,thick] (0,0) -- (1,0);
  \draw[->,thick,blue] (0,0) -- (0.766,0.643) node[above right] {$M$};
  \draw[->] (0.35,0) arc (0:40:0.35);
  \node at (0.52,0.16) {$\theta$};
  \draw[dashed] (0.766,0) -- (0.766,0.643);
  \draw[dashed] (0,0.643) -- (0.766,0.643);
\end{tikzpicture}
\end{center}

% =====================================================================
\clearpage
\section*{Partie 7 — Coordonnées d'un point $M$ du cercle sous la forme $(\cos(\Theta),\sin(\Theta))$}
% =====================================================================

Soit $M$ un point du cercle trigonométrique et soit $\Theta$ la mesure principale de l'angle orienté
\[
\left(\vec{i},\overrightarrow{OM}\right).
\]
Alors
\[
M=\bigl(\cos\Theta,\sin\Theta\bigr).
\]
Autrement dit :
\[
x_M=\cos\Theta,\qquad y_M=\sin\Theta.
\]
On a l'identité fondamentale
\[
\cos^2\Theta+\sin^2\Theta=1.
\]

% =====================================================================
\clearpage
\section*{Partie 8 — Angles et angles associés}
% =====================================================================

\subsection*{8.1. Liaison avec les symétries}

\begin{itemize}
  \item Symétrie d'axe $(O,\vec{i})$ : $\theta\mapsto -\theta$.
  \[
  \cos(-\theta)=\cos\theta,\qquad \sin(-\theta)=-\sin\theta.
  \]

  \item Symétrie d'axe $(O,\vec{j})$ : $\theta\mapsto \pi-\theta$.
  \[
  \cos(\pi-\theta)=-\cos\theta,\qquad \sin(\pi-\theta)=\sin\theta.
  \]

  \item Symétrie de centre $O$ : $\theta\mapsto \pi+\theta$.
  \[
  \cos(\pi+\theta)=-\cos\theta,\qquad \sin(\pi+\theta)=-\sin\theta.
  \]
\end{itemize}

\subsection*{8.2. Liaison avec les rotations}

\begin{itemize}
  \item Rotation d'angle $+\dfrac{\pi}{2}$ : $\theta\mapsto \theta+\dfrac{\pi}{2}$.
  \[
  \cos\left(\theta+\dfrac{\pi}{2}\right)=-\sin\theta,\qquad
  \sin\left(\theta+\dfrac{\pi}{2}\right)=\cos\theta.
  \]

  \item Rotation d'angle $-\dfrac{\pi}{2}$ : $\theta\mapsto \theta-\dfrac{\pi}{2}$.
  \[
  \cos\left(\theta-\dfrac{\pi}{2}\right)=\sin\theta,\qquad
  \sin\left(\theta-\dfrac{\pi}{2}\right)=-\cos\theta.
  \]
\end{itemize}

\subsection*{8.3. Autres angles associés}

\[
\cos\left(\dfrac{\pi}{2}-\theta\right)=\sin\theta,\qquad
\sin\left(\dfrac{\pi}{2}-\theta\right)=\cos\theta.
\]
\[
\cos(2\pi-\theta)=\cos\theta,\qquad
\sin(2\pi-\theta)=-\sin\theta.
\]

\subsection*{8.4. Liaison avec les coordonnées cartésiennes}

Si $M$ a pour coordonnées cartésiennes $(x,y)$ sur le cercle trigonométrique, alors
\[
x=\cos\theta,\qquad y=\sin\theta.
\]
Réciproquement, si $M=(\cos\theta,\sin\theta)$, alors ses coordonnées cartésiennes sont $x=\cos\theta$ et $y=\sin\theta$.

% =====================================================================
\clearpage
\section*{Partie 9 — Tableau récapitulatif}
% =====================================================================

\subsection*{9.1. Symétries}

\begin{center}
\begin{tabular}{|c|c|c|}
\hline
\textbf{Symétrie} & \textbf{Coordonnées de $M'$} & \textbf{Angle image} \\
\hline
Axe $(O,\vec{i})$ & $(x,-y)$ & $-\theta$ \\
\hline
Axe $(O,\vec{j})$ & $(-x,y)$ & $\pi-\theta$ \\
\hline
Centre $O$ & $(-x,-y)$ & $\pi+\theta$ \\
\hline
\end{tabular}
\end{center}

\subsection*{9.2. Rotations}

\begin{center}
\begin{tabular}{|c|c|c|}
\hline
\textbf{Rotation} & \textbf{Coordonnées de $M'$} & \textbf{Angle image} \\
\hline
$+90^\circ=+\dfrac{\pi}{2}$ & $(-y\,;\,x)$ & $\theta+\dfrac{\pi}{2}$ \\
\hline
$-90^\circ=-\dfrac{\pi}{2}$ & $(y\,;\,-x)$ & $\theta-\dfrac{\pi}{2}$ \\
\hline
\end{tabular}
\end{center}

\subsection*{9.3. Angles associés}

\begin{center}
\begin{tabular}{|c|c|c|}
\hline
\textbf{Angle} & $\cos$ & $\sin$ \\
\hline
$\theta$ & $\cos\theta$ & $\sin\theta$ \\
\hline
$-\theta$ & $\cos\theta$ & $-\sin\theta$ \\
\hline
$\pi-\theta$ & $-\cos\theta$ & $\sin\theta$ \\
\hline
$\pi+\theta$ & $-\cos\theta$ & $-\sin\theta$ \\
\hline
$\dfrac{\pi}{2}-\theta$ & $\sin\theta$ & $\cos\theta$ \\
\hline
$\dfrac{\pi}{2}+\theta$ & $-\sin\theta$ & $\cos\theta$ \\
\hline
$2\pi-\theta$ & $\cos\theta$ & $-\sin\theta$ \\
\hline
\end{tabular}
\end{center}

\subsection*{9.4. Angles remarquables}

\begin{center}
\begin{tabular}{|c|c|c|c|c|c|}
\hline
$\theta$ & $0$ & $\dfrac{\pi}{6}$ & $\dfrac{\pi}{4}$ & $\dfrac{\pi}{3}$ & $\dfrac{\pi}{2}$ \\
\hline
Degrés & $0^\circ$ & $30^\circ$ & $45^\circ$ & $60^\circ$ & $90^\circ$ \\
\hline
$\cos\theta$ & $1$ & $\dfrac{\sqrt{3}}{2}$ & $\dfrac{\sqrt{2}}{2}$ & $\dfrac{1}{2}$ & $0$ \\
\hline
$\sin\theta$ & $0$ & $\dfrac{1}{2}$ & $\dfrac{\sqrt{2}}{2}$ & $\dfrac{\sqrt{3}}{2}$ & $1$ \\
\hline
\end{tabular}
\end{center}

\end{document}